\documentclass[11pt]{article}

\usepackage[a4paper,margin=24mm]{geometry}
\usepackage{amsmath,amssymb,amsthm,mathtools}
\usepackage{booktabs}
\usepackage{microtype}
\usepackage{tabularx}
\usepackage[hidelinks]{hyperref}
\usepackage{xcolor}

\newcommand{\trans}{\mathsf{T}}
\newcommand{\E}{\mathcal E}
\newcommand{\K}{\mathcal K}
\newcommand{\Sset}{\mathcal S}
\newcommand{\SLZ}{\mathrm{SL}(2,\mathbb Z)}
\newcommand{\Aut}{\operatorname{Aut}_{\mathbb Z}}
\newcommand{\wordprod}{\Pi}
\newcommand{\Uplus}{U_{+}}
\newcommand{\Uminus}{U_{-}}
\newcommand{\Lplus}{L_{+}}
\newcommand{\Lminus}{L_{-}}

\definecolor{proofblue}{HTML}{1F5A94}
\definecolor{softblue}{HTML}{EDF4FA}
\definecolor{softgray}{HTML}{F3F4F5}

\newtheorem{theorem}{Theorem}
\newtheorem{corollary}{Corollary}
\theoremstyle{definition}
\newtheorem{definition}{Definition}

\title{Why a Shortest Elastic-Reduction Path\\
Contains No Integer Rotation}
\author{Technical note}
\date{14 July 2026}

\begin{document}
\maketitle

\begin{abstract}
The unit-shear elastic reduction searches words in four integer simple
shears.  This note proves that a globally shortest word reaching the elastic
domain cannot contain a nonempty path segment whose product is a lattice
rotation.  The proof uses two exact facts: the elastic domain is invariant
under integer rotations, and conjugating a unit shear by an integer rotation
gives another unit shear.  A rotational segment can therefore be moved to
the end of the path and removed, contradicting shortestness.  This makes the
comparison between the two elastic quadrants precise.  The shortest path is
rotation-free in the path-segment sense, while, for a generic elastic
representative, every path to the other quadrant differs from it by a
quarter turn.  A longer word need not display that quarter turn as a visible
three-shear block; its rotational content is instead exposed by the relative
matrix quotient.
\end{abstract}

\section{Definitions and the exact claim}

Let the four allowed unit-shear generators be
\begin{equation}
 \Sset=\{\Uplus,\Uminus,\Lplus,\Lminus\},
 \qquad
 \Uplus=\begin{pmatrix}1&1\\0&1\end{pmatrix},\quad
 \Lplus=\begin{pmatrix}1&0\\1&1\end{pmatrix},
 \qquad
 \Uminus=\Uplus^{-1},\quad \Lminus=\Lplus^{-1}.
 \label{eq:generators}
\end{equation}
A BFS path is a word \(w=s_1s_2\cdots s_n\) with \(s_i\in\Sset\), and
\(\wordprod(w)=s_1s_2\cdots s_n\in\SLZ\) is its accumulated right-side
transformation.  Starting from a positive-definite metric \(\mathbf C_0\),
the endpoint is
\begin{equation}
 \mathbf C_w=\wordprod(w)^\trans\mathbf C_0\wordprod(w).
 \label{eq:pathaction}
\end{equation}

The integer proper-rotation group is
\begin{equation}
 \K=\mathrm{SO}(2)\cap\SLZ
 =\{\mathbf I,\mathbf Q,-\mathbf I,-\mathbf Q\},
 \qquad
 \mathbf Q=\begin{pmatrix}0&1\\-1&0\end{pmatrix}.
 \label{eq:rotationgroup}
\end{equation}
These are rotations through \(0,-\pi/2,\pi,\pi/2\), respectively.  There
are no others: every column of an integer orthogonal matrix is an integer
unit vector, hence one of the signed coordinate vectors.

\begin{definition}[Rotational segment]
A word \(w\) \emph{contains a nontrivial rotation} if it has a nonempty
contiguous subword \(\rho=s_i\cdots s_j\) for which
\(\wordprod(\rho)\in\K\setminus\{\mathbf I\}\).
\end{definition}

Contiguity is essential: a path segment consists of consecutive moves.  It
would not be meaningful merely to ask whether the final matrix admits some
factorization containing \(\mathbf Q\), because every matrix can be written
trivially as
\(\mathbf M=(\mathbf M\mathbf Q^{-1})\mathbf Q\).

The theoretical elastic domain tested by the implementation can be written
compactly as
\begin{equation}
 \E=\left\{
 \begin{pmatrix}a&b\\b&c\end{pmatrix}\succ0:
 2|b|\leq \min(a,c)
 \right\}.
 \label{eq:elasticdomain}
\end{equation}
Indeed, the four sign-and-swap variants of the fundamental inequalities
\(b\geq0\), \(c\geq a\), and \(2b\leq a\) give exactly
Equation~\eqref{eq:elasticdomain}.  A \emph{globally shortest elastic word}
is a word of minimum length among all words satisfying \(\mathbf C_w\in\E\).
This is the notion of ``short path'' used below.  It minimizes over the full
elastic domain, not over a preselected quadrant.

\section{The two algebraic ingredients}

First, the elastic domain is invariant under every \(\mathbf R\in\K\).  It
is enough to check the generator \(\mathbf Q\):
\begin{equation}
 \mathbf Q^\trans
 \begin{pmatrix}a&b\\b&c\end{pmatrix}
 \mathbf Q
 =\begin{pmatrix}c&-b\\-b&a\end{pmatrix}.
 \label{eq:domaininvariance}
\end{equation}
This swaps \(a\) and \(c\) and changes the sign of \(b\), leaving
\(2|b|\leq\min(a,c)\) unchanged.  Therefore
\begin{equation}
 \mathbf C\in\E
 \quad\Longleftrightarrow\quad
 \mathbf R^\trans\mathbf C\mathbf R\in\E
 \qquad(\mathbf R\in\K).
 \label{eq:Einvariant}
\end{equation}

Second, \(\K\) normalizes the unit-shear generator set.  Direct multiplication
gives
\begin{equation}
 \begin{aligned}
 \mathbf Q\Uplus\mathbf Q^{-1}&=\Lminus,&
 \mathbf Q\Uminus\mathbf Q^{-1}&=\Lplus,\\
 \mathbf Q\Lplus\mathbf Q^{-1}&=\Uminus,&
 \mathbf Q\Lminus\mathbf Q^{-1}&=\Uplus.
 \end{aligned}
 \label{eq:conjugation}
\end{equation}
Conjugation by \(-\mathbf I\) does nothing, and conjugation by
\(-\mathbf Q\) is the same as conjugation by \(\mathbf Q\).  Hence
\begin{equation}
 \mathbf R\Sset\mathbf R^{-1}=\Sset
 \qquad(\mathbf R\in\K).
 \label{eq:normalizes}
\end{equation}
Consequently, if \(v=t_1\cdots t_m\) is a shear word, then
\(\mathbf R\wordprod(v)\mathbf R^{-1}\) is represented by the word obtained
by conjugating every \(t_i\).  The new word has exactly the same length.

\section{The shortest path is rotation-free}

\begin{theorem}[No rotational segment in a shortest elastic word]
Let \(\mathbf C_0\succ0\), and let \(w\) be a globally shortest word in
\(\Sset\) for which \(\mathbf C_w\in\E\).  Then \(w\) contains no nonempty
contiguous subword whose product belongs to
\(\K\setminus\{\mathbf I\}\).
\label{thm:main}
\end{theorem}

\begin{proof}
Assume the opposite.  Factor the word as
\begin{equation}
 w=\alpha\,\rho\,\beta,
 \qquad
 \wordprod(\alpha)=\mathbf A,
 \quad \wordprod(\rho)=\mathbf R\in\K\setminus\{\mathbf I\},
 \quad \wordprod(\beta)=\mathbf B.
\end{equation}
Because \(\mathbf R\) normalizes \(\Sset\), conjugating every letter of
\(\beta\) produces another shear word \(\beta'\) of the same length, with
\begin{equation}
 \wordprod(\beta')=\mathbf R\mathbf B\mathbf R^{-1}.
 \label{eq:conjugatedsuffix}
\end{equation}
The total matrix can now be rearranged exactly:
\begin{equation}
 \wordprod(w)=\mathbf A\mathbf R\mathbf B
 =\mathbf A(\mathbf R\mathbf B\mathbf R^{-1})\mathbf R
 =\underbrace{\mathbf A\wordprod(\beta')}_{\widetilde{\mathbf M}}
  \mathbf R.
 \label{eq:pushrotation}
\end{equation}
Thus the original elastic endpoint is
\begin{equation}
 \mathbf C_w
 =\mathbf R^\trans
  (\widetilde{\mathbf M}^\trans\mathbf C_0\widetilde{\mathbf M})
  \mathbf R\in\E.
\end{equation}
By rotation-invariance of \(\E\), it follows that
\begin{equation}
 \widetilde{\mathbf M}^\trans\mathbf C_0\widetilde{\mathbf M}\in\E.
\end{equation}
But \(\widetilde{\mathbf M}\) is represented by the word
\(\alpha\beta'\), whose length is
\begin{equation}
 |\alpha\beta'|=|w|-|\rho|<|w|.
\end{equation}
This is a strictly shorter path into the elastic domain, contradicting the
minimality of \(w\).  If joining \(\alpha\) to \(\beta'\) creates adjacent
inverse shears, cancelling them only shortens the replacement further.
\end{proof}

\begin{corollary}
A globally shortest elastic word also contains no nonempty contiguous
subword with product \(\mathbf I\).
\end{corollary}

\begin{proof}
Such a subword can be deleted directly without changing the accumulated
matrix.  The resulting word reaches the same endpoint and is shorter.
\end{proof}

\noindent
\colorbox{softblue}{%
\parbox{\dimexpr\textwidth-2\fboxsep}{%
\textbf{What has been proved.}
The short path cannot hide a quarter turn, a half turn, or an identity loop
as any consecutive sequence of its shear moves.  The result is exact and
does not require a generic metric, a unique endpoint, or a trivial lattice
stabilizer.  It relies only on Equations~\eqref{eq:Einvariant} and
\eqref{eq:normalizes}.}}

\section{What this says about the other quadrant}

Theorem~\ref{thm:main} concerns the literal shear word.  The complementary
statement about \emph{all} paths to the two elastic representatives is most
cleanly expressed by a relative matrix quotient.

Let \(\mathbf M_{\mathrm s}\) be a shortest transformation and set
\begin{equation}
 \mathbf C_{\mathrm s}
 =\mathbf M_{\mathrm s}^\trans\mathbf C_0\mathbf M_{\mathrm s}.
\end{equation}
For any other terminal transformation \(\mathbf M\), define
\(\mathbf R=\mathbf M_{\mathrm s}^{-1}\mathbf M\).  Then
\begin{equation}
 \mathbf M^\trans\mathbf C_0\mathbf M
 =\mathbf R^\trans\mathbf C_{\mathrm s}\mathbf R.
 \label{eq:relativeendpoint}
\end{equation}
If the endpoint equals \(\mathbf C_{\mathrm s}\), then
\begin{equation}
 \mathbf R\in\Aut(\mathbf C_{\mathrm s})
 :=\{\mathbf A\in\SLZ:
 \mathbf A^\trans\mathbf C_{\mathrm s}\mathbf A
 =\mathbf C_{\mathrm s}\}.
 \label{eq:stabilizer}
\end{equation}
If instead the endpoint is the quarter-turned representative
\(\mathbf Q^\trans\mathbf C_{\mathrm s}\mathbf Q\), then
\begin{equation}
 \mathbf R=\mathbf A\mathbf Q
 \quad\text{for some}\quad
 \mathbf A\in\Aut(\mathbf C_{\mathrm s}).
 \label{eq:othercoset}
\end{equation}
For a generic positive-definite metric the integer stabilizer is only
\(\{\pm\mathbf I\}\).  Equations~\eqref{eq:stabilizer}--\eqref{eq:othercoset}
then reduce to the sharp classification
\begin{equation}
 \boxed{
 \begin{aligned}
 \text{same representative:}\quad
 &\mathbf M_{\mathrm s}^{-1}\mathbf M\in\{\pm\mathbf I\},\\
 \text{other representative:}\quad
 &\mathbf M_{\mathrm s}^{-1}\mathbf M\in\{\pm\mathbf Q\}.
 \end{aligned}}
 \label{eq:classification}
\end{equation}
This statement includes arbitrarily long words: however many extra shears a
path uses, its net difference from the short path lies in the indicated
rotation class whenever it reaches the same exact representative.

There is an important limitation.  Equation~\eqref{eq:classification} does
\emph{not} imply that every displayed longer word contains a contiguous
three-letter copy of \(\mathbf Q\).  Group relations and cancellations can
spread that rotational content throughout the word.  What is true is:
\begin{enumerate}
 \item the globally shortest word has no removable rotational segment, by
       Theorem~\ref{thm:main}; and
 \item every word reaching the other generic representative has a net
       quarter-turn relative to that shortest word, by
       Equation~\eqref{eq:classification}.
\end{enumerate}
Together these two statements justify assigning the extra lattice rotation
to the other-quadrant description without falsely claiming that it must be
visible as a local block in every possible word.

\section{The simple-shear example}

For
\begin{equation}
 \mathbf F=\begin{pmatrix}1&1.6\\0&1\end{pmatrix},
 \qquad
 \mathbf C_0=\mathbf F^\trans\mathbf F
 =\begin{pmatrix}1&1.6\\1.6&3.56\end{pmatrix},
\end{equation}
the shortest BFS word and the first word reaching the other representative
are
\begin{align}
 w_{\mathrm s}&=\Uminus\Uminus,
 &\mathbf M_{\mathrm s}&=
 \begin{pmatrix}1&-2\\0&1\end{pmatrix},\\
 w_{\ell}&=\Uminus\Lminus\Uplus,
 &\mathbf M_{\ell}&=
 \begin{pmatrix}2&1\\-1&0\end{pmatrix}.
\end{align}
The short word plainly has no rotational subword; Theorem~\ref{thm:main}
shows that this is structural rather than an accident of its length.  The
relative quotient exposes the other path's quarter turn:
\begin{equation}
 \mathbf M_{\mathrm s}^{-1}\mathbf M_{\ell}
 =\mathbf Q,
 \qquad
 \mathbf M_{\ell}=\mathbf M_{\mathrm s}\mathbf Q.
 \label{eq:examplequotient}
\end{equation}
At word level, the same calculation is
\begin{equation}
 (\Uminus\Uminus)^{-1}
 (\Uminus\Lminus\Uplus)
 =\Uplus\Uplus\Uminus\Lminus\Uplus
 \longrightarrow \Uplus\Lminus\Uplus
 =\mathbf Q,
 \label{eq:wordcancellation}
\end{equation}
where the arrow cancels the adjacent \(\Uplus\Uminus\) pair.  The terminal
metrics are therefore the two quarter-turned elastic representatives
\begin{align}
 \mathbf C_{\mathrm s}
 &=\begin{pmatrix}1&-0.4\\-0.4&1.16\end{pmatrix},&
 \mathbf C_{\ell}
 &=\begin{pmatrix}1.16&0.4\\0.4&1\end{pmatrix}
 =\mathbf Q^\trans\mathbf C_{\mathrm s}\mathbf Q.
\end{align}

\begin{table}[htbp]
\centering
\caption{Computational checks used while developing the algebraic statement.
They support the theorem but are not part of its proof.}
\label{tab:checks}
\small
\renewcommand{\arraystretch}{1.22}
\begin{tabularx}{\textwidth}{@{}Xrr@{}}
\toprule
Check & Cases examined & Failures \\
\midrule
Global shortest word had a contiguous nontrivial rotational segment
 & 73 generic starts & 0 \\
All terminal paths obeyed the same/other relative-rotation classification
 & 2414 terminal paths from 75 starts & 0 \\
\bottomrule
\end{tabularx}
\end{table}

\section{Conclusion}

The missing logical step is now supplied: the shortest elastic-reduction
path is not merely shorter than a path carrying a relative quarter turn; it
is provably free of every removable integer-rotation segment.  If such a
segment existed, rotation invariance of the target and rotational closure of
the shear generators would construct a strictly shorter elastic path.

For generic metrics, any path to the other elastic representative differs
from this rotation-free shortest path by \(\pm\mathbf Q\), while paths to the
same representative differ by \(\pm\mathbf I\).  At exceptional symmetric
metrics the stabilizer in Equation~\eqref{eq:stabilizer} can be larger, so
the cosets in Equation~\eqref{eq:classification} must be replaced by the
full stabilizer and its \(\mathbf Q\)-coset.  The rotation-free theorem itself
is unaffected by that exception.

\end{document}
